
\section{问题重述}
\subsection{问题背景}

山区乡村土地资源有限，地块类型与面积差异却十分明显。平旱地、梯田、山坡地、水浇地和大棚适种作物不同，同一作物在不同地块和季次上的亩产量、成本与价格也不同。种植计划若只追逐单季利润，容易破坏跨季轮作和豆类养地要求，还可能产生过度集中或难以管理的碎片化安排。

本题要求给出2024至2030年的完整种植面积表。第一问参数保持稳定，但需要分别处理滞销与折价销售。第二问加入销量、亩产量、成本和价格的不确定性。第三问进一步要求考虑作物间关系及经济变量相关性。已有研究表明，数学规划能够系统表达土地和轮作约束\upcite{tan2025review,xu2025lp,dantzig1963}，尾部风险指标可补充平均收益评价\upcite{shi2025cvar,rockafellar2000}。不过，附件只有2023年单年截面，无法直接估计未来概率分布与替代弹性，因此方法必须把观测数据和模拟假设分开。

本文采用分阶段建模。先在稳定参数下建立可验证的多期混合整数线性规划，再用中心趋势求得固定方案并进行样本外风险评估，最后在受控模拟中加入关系结构。这样既能输出逐年逐季逐地块的完整面积方案，也能说明收益优势在哪些情景下成立，哪些结论不能外推。


\subsection{问题重述}

\textbf{问题1} 在2023年亩产量、成本、价格和销量代理保持不变的条件下，分别针对超额产量滞销与按正常价格50\%销售两种情况，制定2024至2030年的种植方案并填写结果表。

\textbf{问题2} 在预期销量、亩产量、成本和价格随年份波动的条件下，给出一套兼顾高期望收益与下行风险的种植方案，并填写结果表。

\textbf{问题3} 在问题二基础上加入作物替代、互补关系及经济变量相关性，比较关系结构对收益分布和种植面积配置的影响。
